% Original applied explanation; see Chapter04/SOURCES.md.
\subsubsection{ট্যাবের নাম ও পেজের শিরোনাম কেন আলাদা?}
কলেজের Notice পেজে \texttt{<title>} ব্রাউজারের ট্যাবের পরিচয় দেয়; \texttt{<h1>} পেজের ভেতরে শিরোনাম দেখায়। \texttt{head}-এ নথি সম্পর্কে তথ্য থাকে, আর \texttt{body}-তে দর্শকের জন্য মূল বিষয়বস্তু থাকে। \texttt{<!DOCTYPE html>} আধুনিক HTML নথি হিসেবে কাজ করার ঘোষণা।

বাংলা লেখার জন্য UTF-8 ঘোষণা ও ফাইলের প্রকৃত UTF-8 সংরক্ষণ মিলতে হবে। ভুল অক্ষর দেখলে শুধু ট্যাগ বদলাবে না; ফাইলের encoding-ও দেখবে। সম্পাদনার পরে ফাইল সংরক্ষণ করে ব্রাউজার আবার লোড করলে নতুন সংস্করণ পরীক্ষা করা যায়।
